\documentclass[11pt]{article}
\usepackage[margin=2.2cm]{geometry}
\usepackage{amsmath,amssymb,amsthm,booktabs,enumitem}
\newtheorem{theorem}{Theorem}[section]
\newtheorem{lemma}[theorem]{Lemma}
\newtheorem{corollary}[theorem]{Corollary}
\newtheorem{proposition}[theorem]{Proposition}
\newtheorem{conjecture}[theorem]{Conjecture}
\theoremstyle{definition}
\newtheorem{definition}[theorem]{Definition}
\theoremstyle{remark}
\newtheorem{remark}[theorem]{Remark}
\newcommand{\hG}{h_\Gamma}
\newcommand{\Gh}{\Gamma_h}
\newcommand{\Oh}{\Omega_h}
\newcommand{\Zh}{Z_h}
\newcommand{\Yh}{Y_h}
\newcommand{\norm}[1]{\|#1\|}
\newcommand{\tnorm}[1]{|\!|\!|#1|\!|\!|}
\newcommand{\ip}[2]{\langle #1,#2\rangle}
\newcommand{\dn}{\partial_n}
\newcommand{\Ut}{\tilde U}
\newcommand{\ut}{\tilde u}
\newcommand{\pbar}{\bar p}
\newcommand{\GS}{\mathrm{GS}}
\newcommand{\CNS}{\mathrm{CNS}^\ast}
\newcommand{\Re}{\mathrm{Re}}
\newcommand{\lab}[1]{\textbf{[#1]}}
\renewcommand{\div}{\operatorname{div}}
\title{notes/v7/nsconst: the Navier--Stokes constants $C_R,\beta_O$ and the explicit $h_0$ for the paper's test;\\
the limit $\gamma^\ast_\infty$ of the coercivity threshold; convergence of the leak constant for $\gamma\hG\to\infty$}
\date{2026-09-28}
\begin{document}\maketitle

\section*{0. Summary}
Nothing in \texttt{paper/}, \texttt{paper3d/}, \texttt{letter/} was edited; nothing was committed. Scripts and logs are in this directory (Section~\ref{sec:files}); every run is $<1$\,min on one core.

\begin{enumerate}[label=(\alph*)]
\item \lab{NUMERICAL} \emph{Constants of Theorem~\texttt{n2:L} for the paper's Navier--Stokes test} (Section~\ref{sec:consts}). The only NS test with data is test~B of \texttt{code/job\_v4.py ns} (\texttt{results/v4/report\_v4.txt}, part~B; there is no NS test in \texttt{13\_numerics.tex}): $\Omega=(-2.5,2.5)^2\setminus\overline{B_1}$, $u=\operatorname{curl}\bigl((r^2-1)^2(x+y^2/2)/10\bigr)$, $\nu\in\{1,0.1\}$. Computed on the \emph{exact} domain with an isoparametric Taylor--Hood discretisation (validated), with convergence tables:
\[
U_\infty=175.6,\quad C_F=0.56423,\quad \beta_O=0.8225\ (\nu=1),\ \ 0.0533\ (\nu=0.1),\quad b_O=\beta_O/\nu=0.823,\ 0.533,
\]
\[
C_R^{\rm quad}\ge2.477\ (\text{converged lower-bound sequence}),\qquad C_R\le\sqrt2\,C_R^{\rm quad}\approx3.50\ \text{(sum form of (R$_S$))}.
\]
The maximiser of $C_R$ is a global, pressure-dominated mode (Neumann--Poincar\'e); the flat-boundary essential part is only $1.538$ (Lemma~\ref{lem:flat}).
\item \lab{NUMERICAL, negative} \emph{The explicit $h_0$.} The sufficient condition \eqref{n2:h0} becomes
$(1+\Re)h+(\hG/\gamma_0)^{1/2}\le c_\natural\cdot7.6\cdot10^{-6}$ ($\nu=1$, $\Re=176$) and $\le c_\natural\cdot4.9\cdot10^{-8}$ ($\nu=0.1$, $\Re=1756$),
i.e.\ $h\le4.3\cdot10^{-8}c_\natural$, resp.\ $2.8\cdot10^{-11}c_\natural$, with $c_\natural$ (through $C_\sharp$) still not computed.
The production meshes have $h=0.11$--$1.5$. The theorem is therefore qualitative for this test. The \emph{actual} discrete constant of (L) on the production meshes is $\beta_L/\nu=0.021,0.016,0.014$ ($N=16,32,48$, $\mu=100$), equal to the Stokes value to $0.1$--$1\%$ for both $\nu$: the Oseen perturbation is invisible at $N=16$, and what degrades $\beta_L$ is the closeness of $\mu=100$ to the coercivity threshold.
\item \lab{PROVED MODULO (CUT), (LOC)} \emph{$\gamma^\ast_\infty$} (Theorem~\ref{thm:gam}). On the production family, $\liminf_N\gamma^\ast(N)\ge\Gamma_\infty$ (modulo a standard spline cutoff lemma (CUT)) and $\limsup_N\gamma^\ast(N)\le\Gamma_\infty$ (modulo a discrete localisation inequality (LOC)), where
$\Gamma_\infty:=\sup_{|\theta|\le\pi/4}\gamma_c(a(\theta))/\cos^2\theta$, $a(\theta)=(L\sec\theta-1)/(2\cos^2\theta)$, and $\gamma_c(a)$ is the half-plane lattice threshold. The Floquet reduction of $\gamma_c$ is PROVED (Lemma~\ref{lem:floq}).
\lab{NUMERICAL}: the supremum is at $\theta=0$ and Floquet phase $0$, and
\[
\Gamma_\infty=\gamma_c(\tfrac34)=21.4080\ \ (\text{all depths}),\qquad 21.4045\ (2\text{ layers}),\qquad 21.3860\ (1\text{ layer}),
\]
which confirms the paper's extrapolated $21.4$ (window Richardson: $21.39$ / $21.41$ for 1 / 2 layers). The two-layer global thresholds for $N=128,\dots,2048$ (new) extrapolate to $21.40452$, against the cell value $21.40452$; the $1/N$ coefficient $21.8$ is predicted as $22.0$ by a harmonic-well/curvature expansion.
\item \lab{PROVED} \emph{Leak constant for $\gamma\hG\to\infty$} (Corollary~\ref{cor:K}). The item ``numerical only'' of \texttt{14\_open} l.26 (first half) is already a consequence of the paper's Theorem~\texttt{u:up} and Corollary~\texttt{u:K}: on test~P, $K_h\to K$ for \emph{every} family with $\mu\ge\mu_0$, $|K_h-K|\le C(\hG+(E^0_U)^2)$, with no relation between $\gamma$ and $\hG$. For general data the normalised $\GS$ $H^1$ constant has a PROVED dichotomy: $\to K$ if $\gamma\hG^{1/2}\to0$, $\to\infty$ if $\gamma\hG^{1/2}\to\infty$ (in particular if $\gamma\hG\to\infty$) whenever $\dn u\cdot\tau\not\equiv0$. Navier--Stokes analogue: $K_{NS,h}\to K_{NS}$ for every $\gamma\ge\gamma_0$, rate $\hG^{1/2}$, PROVED MODULO (R$_S$).
\item \lab{CONJECTURE + NUMERICAL} \emph{$\lim d_h/\hG^{1/2}$ on LP meshes} (Conjecture~\ref{conj:dh}). The cell problem gives an explicit candidate $D(\mu)$; for the production meshes and test~P it predicts $D(10^8)=0.564$ and $D(100)=0.2256$, against the fitted $0.55$ and $0.225$ of \texttt{u:rem}. A proof needs a transfer lemma for $C^1$ quintic splines between $O(\hG)$-close lattices and exponential decay of the cell corrector; it is not given.
\end{enumerate}

\section{The constants of Theorem \texttt{n2:L} for the paper's test}\label{sec:consts}

\subsection{What is to be computed}
In \texttt{notes/v6/ns2} (paper Section~\texttt{n2}): $U_\infty=\max(\norm{\ut}_{L^\infty(\Omega\cup A)},\norm{\nabla\ut}_{L^\infty(\Omega\cup A)})$;
$C_F$ the Friedrichs constant; $\beta_O=\inf_{w\in V}\sup_{v\in V}\langle Lw,v\rangle/(\norm{\nabla w}\norm{\nabla v})$ with
$Lw=-\nu\Delta w+(w\cdot\nabla)u+(u\cdot\nabla)w$ on $V=\{v\in H^1_0(\Omega)^2:\div v=0\}$; $C_R$ the constant of (R$_S$)(a):
$\norm\varphi_{H^2(\Omega)}+\norm{\pi-\bar\pi}_{H^1(\Omega)}\le C_R\norm g$ for $-\Delta\varphi+\nabla\pi=g$, $\varphi\in V$. We write
\[
C_R^{\rm quad}:=\sup_g\Bigl(\frac{\norm\varphi^2_{H^2}+\norm{\pi-\bar\pi}^2_{H^1}}{\norm g^2}\Bigr)^{1/2},\qquad C_R^{\rm quad}\le C_R\le\sqrt2\,C_R^{\rm quad},
\]
with full norms $\norm\varphi^2_{H^2}=\norm\varphi^2+\norm{\nabla\varphi}^2+\norm{D^2\varphi}^2$.

\subsection{Method}
\texttt{iso\_th.py}: Taylor--Hood $P_k/P_{k-1}$ (continuous pressure) on the \emph{exact} $\Omega$, production-mesh topology (\texttt{svn.make\_mesh}), boundary triangles curved by $P_k$-isoparametric interpolation of the Zl\'amal blending of the radial projection, exact Jacobians at all quadrature points. Validation (\texttt{consts.py check}): area error $3.4\cdot10^{-5}$, $9.2\cdot10^{-7}$ ($N=8,16$; rate $h^5$); manufactured Stokes solution, relative $H^1$ error $0.274,0.056,0.0061$ ($N=8,16,32$); $\beta_O=\nu$ to $10^{-15}$ for $u=0$.

$\beta_O$: $1/\beta_O^2$ is the largest eigenvalue of $y\mapsto L^{-\mathsf T}AL^{-1}Ay$ on the discrete $V_h$ ($A$ the vector Dirichlet stiffness, $L^{-1}$ a saddle solve); ARPACK.

$C_R$: \emph{Grisvard's identity} avoids second derivatives of the discrete solution. For $w\in H^2\cap H^1_0$ on a curvilinear polygon with convex corners, $\norm{D^2w}^2=\norm{\Delta w}^2-\int_{\partial\Omega}\kappa|\dn w|^2$ \cite[Thm~3.1.1.1, Lemma~4.3.1.3]{Grisvard}; here $\kappa=0$ on the square and $\kappa=-1$ on $\Gamma$ (concave for $\Omega$), and $\Delta\varphi=\nabla\pi-g$. Hence
\[
\norm\varphi^2_{H^2}+\norm{\pi-\bar\pi}^2_{H^1}=\norm\varphi^2+\norm{\nabla\varphi}^2+\norm{\nabla\pi-g}^2+\norm{\dn\varphi}^2_{\Gamma}+\norm{\pi-\bar\pi}^2+\norm{\nabla\pi}^2 ,
\]
which needs only first derivatives (the Taylor--Hood pressure is continuous). The supremum is taken over $g\in\mathcal P_M:=\{$vector polynomials of total degree $\le M$ on the box, restricted to $\Omega\}$, with a fine solve for each basis function. Since $\mathcal P_M\subset\mathcal P_{M+1}$ and $\bigcup_M\mathcal P_M$ is dense in $L^2(\Omega)^2$, $\sup_{\mathcal P_M}\uparrow(C_R^{\rm quad})^2$ as $M\to\infty$ (for exact solves). For each $M$ the fine solves are converged in $N$.

\subsection{Results}
$U_\infty$ (\texttt{consts.py Uinf}): $\max_{\overline\Omega}|u|=125.0$, $\max_{\overline\Omega}|\nabla u|_F=175.6$ (operator norm $173.3$), both at the corners of the box; on the collar $0.9<r<1$ the values are $0.06$, $0.8$. So $U_\infty=175.6$ and $\Re=U_\infty/\nu=175.6$, $1756$.

Friedrichs (\texttt{consts.py CF}, $N=16,32$ agree to $2\cdot10^{-6}$): $\lambda_1^{\rm Dir}(\Omega)=3.14111$, $C_F=0.564233$ for $H^1_0(\Omega)$ (used in Lemma~\texttt{n2:reg}); $C_F^R=0.906591$ for fields vanishing on the square only (Neumann on $\Gamma$; the constant of Lemma~\texttt{lem:korn} in the limit).

\begin{table}[h]\centering\small
\caption{$\beta_O$ (\texttt{run\_beta.log}).}\label{tab:beta}
\begin{tabular}{l|cccc|cc|c}
\toprule
 & \multicolumn{4}{c|}{$k=4$} & \multicolumn{2}{c|}{$k=6$} & \\
$\nu$ & $N=8$ & 16 & 32 & 48 & $N=16$ & 24 & estimate\\
\midrule
1 & .82205 & .82109 & .82238 & .82253 & .82281 & .82256 & $0.8225\pm0.0002$\\
0.1 & .04771 & .05469 & .05357 & .05335 & .05441 & .05364 & $0.0533\pm0.0002$\\
\bottomrule
\end{tabular}
\end{table}

\begin{table}[h]\centering\small
\caption{$C_R^{\rm quad}$ lower-bound sequence (\texttt{run\_CR.log}, \texttt{run\_CR16.log}), and the value of each part of the form at the maximiser ($N=32$, $M=16$).}\label{tab:CR}
\begin{tabular}{r|ccccc}
\toprule
$M$ & 4 & 8 & 12 & 16 & 20\\
\midrule
$N=16$ & 2.3477 & 2.4595 & 2.4710 & 2.4755 & 2.4768\\
$N=32$ & -- & 2.4588 & 2.4704 & 2.4748 & --\\
\bottomrule
\end{tabular}\\[3pt]
$\norm\varphi^2:.0003$, $\norm{\nabla\varphi}^2:.0068$, $\norm{\Delta\varphi}^2:.362$, $\norm{\dn\varphi}^2_\Gamma:.038$, $\norm{\pi-\bar\pi}^2:4.485$, $\norm{\nabla\pi}^2:1.232$ (times $\norm g^2$).
\end{table}
The sequence increases by $0.011$, $0.0045$, $0.0013$ per step in $M$; a geometric extrapolation gives $C_R^{\rm quad}\approx2.478$. The maximiser is almost a gradient, $g\approx\nabla q$ with $q$ the first nonconstant Neumann eigenfunction of $\Omega$: then $\varphi\approx0$, $\pi\approx q$ and the value is $1+1/\mu_2^N$ ($\mu_2^N\approx0.275$).

\begin{lemma}[Flat-boundary part; NUMERICAL]\label{lem:flat}
For the half-plane and one Fourier mode (every mode gives the same value by scaling),
$q_{\rm flat}:=\sup_g(\norm{D^2\varphi}^2+\norm{\nabla\pi}^2)/\norm g^2=2.366$ (\texttt{halfplane.py}; $2.3665,2.3671,2.3660$ on three resolutions; numerically $(3+\sqrt3)/2$). In the interior the value is $1$ ($D^2\varphi\leftrightarrow Pg$, $\nabla\pi\leftrightarrow(I-P)g$). So the essential spectrum coming from the flat sides and from $\Gamma$ lies at $\le1.538^2$, below the computed $2.478^2$.
\end{lemma}
The corner contribution (quarter-plane problem) was \emph{not} computed; it is the one gap in the claim ``$C_R^{\rm quad}=2.478$''. The quantity entering \eqref{n2:h0} is the sum form: the minimax bound $(a+b)^2\le(1+s)a^2+(1+1/s)b^2$, minimised over $s$, gives $3.16$ ($M=16$), $3.39$ ($M=20$) and still grows with $M$ (the velocity part approaches its essential value); we use $C_R=\sqrt2\cdot2.478=3.50$. (\texttt{run\_CR.log} predates the minimax; its column \texttt{CR\_sum\_bound} is the $\sqrt2$ bound.)

\subsection{The threshold $h_0$}
With these values, Lemma~\texttt{n2:reg} gives $\Lambda_O=1+(1+C_F)C_FU_\infty/\beta_O=189.4$ ($\nu=1$), $2909$ ($\nu=0.1$), so $\Theta_R=\max(1,C_R\Lambda_O)=663$, resp.\ $1.02\cdot10^4$: the dual solution of Lemma~\texttt{n2:dual} is controlled only with a factor $10^3$--$10^4$. In the form of Remark~\texttt{n2:rem:h0},
\[
(1+\Re)h+\Bigl(\frac\hG{\gamma_0}\Bigr)^{1/2}\le\frac{c_\natural}{\max(1,C_R)\,\Re\,(1+\Re/b_O)}=
\begin{cases}c_\natural/(3.50\cdot175.6\cdot214.5)=7.6\cdot10^{-6}\,c_\natural,&\nu=1,\\ c_\natural\,0.1/(3.50\cdot1756\cdot3296)=4.9\cdot10^{-8}\,c_\natural,&\nu=0.1.\end{cases}
\]
Hence $h_0\le4.3\cdot10^{-8}c_\natural$ and $\hG\le5.8\cdot10^{-11}\gamma_0c_\natural^2$ ($\nu=1$); $h_0\le2.8\cdot10^{-11}c_\natural$ and $\hG\le2.4\cdot10^{-15}\gamma_0c_\natural^2$ ($\nu=0.1$).
\lab{NOT COMPUTED}: $c_\natural$, i.e.\ $C_\sharp$ (a product of about fifteen trace, inverse, extension and interpolation constants of Lemmas~\texttt{n2:Y}, \texttt{n2:dual}). Its value cannot rescue the estimate: at the production meshes the left side is $\ge20$ ($\nu=1$), so one would need $c_\natural\gtrsim3\cdot10^6$.

\paragraph{Where the pessimism comes from.} The factor $\Re(1+\Re/b_O)\approx\Re^2/b_O$ comes from bounding $K$ by $U_\infty$ (the corner maximum of $|\nabla u|$, $175$; the field is $\le1$ near $\Gamma$) and from the duality chain. The stability margin itself is harmless ($b_O=0.82,0.53$). The actual discrete constant (\texttt{discrete\_L.py}, $\GS$, $\mu=100$, energy norm):
\begin{center}\small
\begin{tabular}{r|ccc}
\toprule
$N$ & 16 & 32 & 48\\
\midrule
$\beta_L/\nu$, $\nu=1$ & .02125 & .01623 & .01421\\
$\beta_L/\nu$, $\nu=0.1$ & .02140 & .01629 & --\\
Stokes ($K=0$) & .02122 & .01622 & .01420\\
\bottomrule
\end{tabular}
\end{center}
(L) holds on every production mesh, with the Stokes constant. The decrease with $N$ tracks the margin $1-\mu^\ast(N)/\mu$ ($0.41,0.34,0.31$), cf.\ Section~\ref{sec:gam}. So for this test the explicit $h_0$ is off by at least seven orders of magnitude. The $\Re^2/b_O$ scaling is the known worst case of Schatz's argument, and the discrete (L) is observed from $N=16$.

\paragraph{Status.} (R$_S$) is a qualitative theorem (Dauge); the numbers above make it quantitative for this $\Omega$ up to the corner gap of Lemma~\ref{lem:flat}. They are NUMERICAL (discretisation plus extrapolation, with the convergence tables above), not rigorous enclosures. A rigorous $\beta_O$ would need a guaranteed lower bound on the smallest singular value (e.g.\ a Lehmann--Goerisch or residual argument on top of Table~\ref{tab:beta}); this was not attempted.

\section{The limit of the coercivity threshold}\label{sec:gam}

\subsection{Setting}
$\gamma^\ast(N)=\hG\sup_{v\in\Zh}(2B-A)/C$ on the production mesh $\mathcal T_N$ (\texttt{svn.make\_mesh}, $L=2.5$, $k=4$), with $A(v)=a_h(v,v)$, $B(v)=\ip{\dn v}{v}_{\Gh}$, $C(v)=\norm v^2_{\Gh}$, $\hG=\max|e|$ (paper, Section~\texttt{pb:sec}).

\emph{Local geometry.} On the side $x=L$ the wall vertex at polar angle $\theta$ has chord $\ell(\theta)=\hG\cos^2\theta\,(1+O(N^{-1}))$ and the layers along each ray have thickness $(L\sec\theta-1)/m$, $m=N/4$. After scaling by the chord, the mesh in a window of $W$ columns and $D$ layers around $\theta$ converges, with vertex errors $O((W+D)/N)$, to the rectangular lattice
\[
\mathcal L_a:\ \text{columns }[i,i+1],\ \text{rows }[ja,(j+1)a],\ \text{forward diagonals},\qquad a=a(\theta)=\frac{L\sec\theta-1}{2\cos^2\theta},
\]
with the wall at the bottom row (checked in \texttt{cell\_limit.py aspect}: at $N=256$ the measured chord ratio equals $\cos^2\theta$ to $10^{-4}$ and the aspect equals $a(\theta)$ to $1\%$). The other sides follow by the rotation symmetry of the mesh; it is not a reflection symmetry, but $\mathcal L_{a(\theta)}$ and $\mathcal L_{a(-\theta)}$ coincide. Local penalty in edge units: $(\mu/h)\ell(\theta)=\gamma\cos^2\theta$.

\begin{definition}
$Z(\mathcal L_a)$: continuous piecewise $P_4$, exactly divergence-free fields with compact support in $\overline{\mathcal L_a}$ (no condition on the wall). $\gamma_c(a):=\sup\{(2B-A)/C:\ v\in Z(\mathcal L_a),\ C(v)>0\}$ (edge units). For $J\in\mathbb N$ and a phase $\phi$, $\gamma_c(a,J,\phi)$ is the same supremum over $\phi$-quasi-periodic fields on one column of $J$ rows, zero on row $J$ (the one-edge cell of \texttt{unifmu2/floquet\_cell.py} with $N\to\infty$).
\end{definition}

\begin{lemma}[Floquet reduction; PROVED for ``$\le$'', PROVED MODULO (CUT) for ``$\ge$'']\label{lem:floq}
$\gamma_c(a)=\sup_J\sup_\phi\gamma_c(a,J,\phi)$, and $J\mapsto\sup_\phi\gamma_c(a,J,\phi)$ is nondecreasing.
\end{lemma}
\begin{proof}
($\le$) Let $v\in Z(\mathcal L_a)$ be supported in $W$ columns and $J$ rows, and let $W'\ge W$. The $W'$-periodisation $\sum_{m\in\mathbb Z}v(\cdot-mW'e_1)$ is continuous (the translates meet only where they vanish) and divergence-free, and its forms per period equal those of $v$. On the $W'$-periodic supercell the forms commute with the translation by one column, so the discrete Fourier transform over $\mathbb Z_{W'}$ splits the supercell pencil $(2B-A,C)$ orthogonally into the $W'$ one-column pencils with phases $2\pi j/W'$. Hence $(2B-A)(v)/C(v)\le\max_j\gamma_c(a,J,2\pi j/W')$. Monotonicity in $J$: extend by zero.
($\ge$) Let $w$ be a (complex) Bloch witness with phase $\phi$; since the forms have real coefficients, $\operatorname{Re}w$ or $\operatorname{Im}w$ has at least the same quotient over one period. Its stream function $\psi$ is a $C^1$ piecewise quintic, vanishing with its gradient on row $J$, and $e^{i\phi}$-quasi-periodic in $x$: the flux through a vertical segment changes over one period by the wall flux of one cell, which is $0$. (CUT) gives $\psi_W\in S^1_5$ equal to $\psi$ on $W$ columns, zero outside $W+2K$ columns, with $\norm{D^2\psi_W}^2+\norm{\nabla\psi_W}^2_{\rm wall}\le C_w$ on the $2K$ transition columns. Then $v_W=\operatorname{curl}\psi_W\in Z(\mathcal L_a)$ has all three forms equal to $W$ times the one-period values up to $O(1)$, and the quotient tends to that of $w$.
\end{proof}
\begin{itemize}
\item[(CUT)] \emph{Cutoff in $S^1_5(\mathcal L_a)$}: for some $K$ and every $s\in S^1_5$ on a band of $K$ columns there is $\tilde s\in S^1_5$ that agrees with $s$ near the left edge of the band, vanishes near its right edge, and satisfies $\norm{\tilde s}_{H^2(\text{band})}\le C\norm s_{H^2(\text{band})}+C\norm s_{L^2}$.
\end{itemize}
(CUT) follows from the existence of stable local bases (stable local minimal determining sets) for $C^1$ splines of degree $\ge5$ on arbitrary triangulations \cite[Ch.~7--9]{LaiSchumaker} %% VERIFY: exact theorem number for stable local MDS of S^1_5
: take $\tilde s=\sum_j\chi(x_j)c_j(s)B_j$ with a $0/1$ step $\chi$. The lattice has no singular vertices (every interior vertex lies on three lines).

\begin{theorem}[$\gamma^\ast_\infty$]\label{thm:gam}
Let $\Gamma_\infty:=\sup_{0\le\theta\le\pi/4}\gamma_c(a(\theta))/\cos^2\theta$.
\begin{enumerate}[label=(\alph*)]
\item \lab{PROVED MODULO (CUT)} $\liminf_{N\to\infty}\gamma^\ast(N)\ge\Gamma_\infty$.
\item \lab{PROVED MODULO (LOC)} $\limsup_{N\to\infty}\gamma^\ast(N)\le\Gamma_\infty$.
\end{enumerate}
Hence, modulo (CUT) and (LOC), $\gamma^\ast_\infty=\lim_N\gamma^\ast(N)$ exists and equals $\Gamma_\infty$.
\end{theorem}
\begin{proof}
(a) Fix $\theta_0$ and $\varepsilon>0$. By Lemma~\ref{lem:floq} there is $v\in Z(\mathcal L_{a(\theta_0)})$, supported in a window $P$ of $W$ columns and $J$ rows, with quotient $>\gamma_c(a(\theta_0))-\varepsilon$. Let $P_N$ be the corresponding patch of $\mathcal T_N$ at $\theta_0$, scaled by the chord. Its vertices converge to those of $P$ and no vertex of $P$ is singular. So the dimension of the space of $P_4$ divergence-free fields supported in $P_N$ (the Morgan--Scott count for $S^1_5$ with its boundary conditions) is constant for large $N$, and the space depends continuously on the vertices. This is the perturbation argument of Proposition~\texttt{prop:penalty}(b), used there for $\cS_{3/4}$ and the $W=4$ window. So the patch quotient converges to $(2B-A)(v)/C(v)$. Undoing the scaling, a field supported in $P_N$ gives $\gamma^\ast(N)\ge(\hG/\ell(\theta_0))\cdot$ (patch quotient) $\to(\gamma_c(a(\theta_0))-\varepsilon)/\cos^2\theta_0$.

(b) Let $\gamma>\Gamma_\infty$ and $\lambda=\gamma/\hG$. We must show $N_h=A-2B+\lambda C\ge0$ on $\Zh$ for large $N$. Three ingredients:
\begin{itemize}
\item \emph{Window bound} (PROVED). For fixed $W,D$, uniformly over the windows along $\Gamma$, every $v\in\Zh$ supported in a window of $W$ columns and $D$ layers at angle $\theta$ satisfies $2B-A\le(\gamma_c(a(\theta))/\cos^2\theta+o(1))\,C/\hG$. Proof: the patch geometry converges uniformly in $\theta$ (compactness of $[0,\pi/4]$, Lipschitz dependence of the local lattice on $\theta$), the patch quotient is continuous at nonsingular configurations, and a window of the lattice is bounded by $\gamma_c$ by the ``$\le$'' half of Lemma~\ref{lem:floq}. The corners $\theta=\pi/4$ need no special treatment: $\ell$ and $a$ are continuous there and only their derivatives jump.
\item \emph{Coercivity at a high level} (PROVED, paper Lemma~\texttt{lem:coercive}): $A-2B+\lambda_0C\ge c_1(A+\lambda_0C)$ for $\lambda_0=\gamma_0/\hG$.
\item \emph{(LOC)}, below.
\end{itemize}
Given (LOC) with $\varepsilon$: $A-2B+\lambda C\ge\sum_i(A-2B+\lambda C)(v_i)-\varepsilon(A+\hG^{-1}C)(v)\ge-\varepsilon(A+\hG^{-1}C)(v)$ once $\lambda\hG>\Gamma_\infty+o(1)$. Take a convex combination with the coercive level $\lambda_0$: for $t=\varepsilon/(c_1+\varepsilon)$ and $\lambda'=(1-t)\lambda+t\lambda_0$, $N_h(\lambda')(v,v)\ge0$. Since $\lambda'\hG\to\gamma$ as $\varepsilon\to0$, the claim follows.
\end{proof}

\begin{itemize}
\item[(LOC)] \emph{Discrete IMS localisation.} For every $\varepsilon>0$ there are $W,D,N_0$ such that for $N\ge N_0$ every $v\in\Zh$ can be written as $v=v_0+\sum_iv_i$, where $v_i\in\Zh$ is supported in the $i$-th boundary window ($W$ columns, $D$ layers, windows overlapping at most twice), $v_0\in\Zh\cap V_O$, and
$(A-2B+\lambda C)(v)\ge\sum_{i\ge1}(A-2B+\lambda C)(v_i)-\varepsilon\,(A(v)+\hG^{-1}C(v))$ for $\lambda\hG$ in a compact set.
\end{itemize}
\emph{Why (LOC) is plausible, and what a proof needs.} For $H^1$ fields and smooth partitions $\sum\chi_i^2=1$ the IMS formula is exact for the boundary forms ($\sum_iB(\chi_iv)=B(v)$ because $\sum\chi_i\dn\chi_i=0$) and costs $\sum\norm{\nabla\chi_i\otimes v}^2$ in $A$. Along $\Gamma$ this is $O((D/W)^2)A+O(D/W^2)\hG^{-1}C$ (1D Poincar\'e from the wall). In depth, a logarithmic cutoff over layers $D\dots D^2$ with Hardy's inequality $\int|v-v(0)|^2/y^2\le4\int|\partial_yv|^2$ gives $O(1/\log^2D)$. The missing step is to put $\chi_iv$ back into $\Zh$. The natural route is the stream function and a local quasi-interpolant onto $S^1_5$ that reproduces quintics (it exists by the same local-basis theory as (CUT)). One must then control lower-order terms of $\psi$ after subtracting local linear polynomials, and keep the wall trace. This is routine but long, and it is not written out.

\subsection{Numerics \lab{NUMERICAL}}\label{sec:gamnum}
\texttt{cell\_limit.py} (reuses \texttt{floquet\_cell.Cell} unchanged).
\begin{itemize}
\item \emph{Phase.} For $a=\frac34$, $J=2$, $N=4096$: $\gamma_c(\phi)$ decreases monotonically from $21.4066$ ($\phi=0$) to $14.917$ ($\phi=\pi$) (\texttt{run\_table.log}, phase scan in the text of the log header). The maximiser is the periodic mode, for every $a$ in the $\theta$ scan.
\item \emph{Depth and flatness} (\texttt{run\_table.log}):
\begin{center}\small
\begin{tabular}{r|cccccc}
\toprule
$J$ & 1 & 2 & 3 & 4 & 6 & 8\\
\midrule
$N=1024$ & 21.3951 & 21.4128 & 21.4161 & 21.4161 & 21.4161 & 21.4161\\
$N=4096$ & 21.3883 & 21.4066 & 21.4100 & 21.4100 & 21.4100 & 21.4100\\
$N=16384$ & 21.3866 & 21.4050 & 21.4085 & 21.4085 & 21.4085 & 21.4085\\
\bottomrule
\end{tabular}
\end{center}
The curvature correction is exactly $O(1/N)$ (the differences shrink by $4.0$); extrapolated, $\gamma_c(\frac34)=21.4080$ ($J\ge3$), $21.4045$ ($J=2$), $21.3860$ ($J=1$). The paper's window Richardson values $21.39$ / $21.41$ (one / two layers) agree.
\item \emph{Which $\theta$} (\texttt{run\_theta.log}, $J=4$): $f(\theta)=\gamma_c(a(\theta))/\cos^2\theta$ decreases monotonically from $21.410$ ($\theta=0$) to $12.61$ ($\theta=\pi/4$; $a=2.54$, $\gamma_c=6.31$). Hence $\Gamma_\infty=\gamma_c(\frac34)=21.408$.
\item \emph{Global thresholds at large $N$} (new; fields supported in the first two layers, which differ from the global threshold by $\le0.002$ for $N\le64$; \texttt{run\_mesh.log}):
\begin{center}\small
\begin{tabular}{r|ccccccc}
\toprule
$N$ & 32 & 64 & 128 & 256 & 512 & 1024 & 2048\\
\midrule
$\gamma^\ast_{J=2}$ & 19.8505 & 20.8024 & 21.1654 & 21.3019 & 21.3576 & 21.3822 & 21.3936\\
$N(\Gamma^{J=2}_\infty-\gamma^\ast_{J=2})$ & -- & -- & 30.6 & 26.3 & 24.0 & 22.9 & 22.4\\
\bottomrule
\end{tabular}
\end{center}
A fit $\Gamma-a/N-b/N^2$ on $N=512$--$2048$ gives $\Gamma=21.40452$, and on $256$--$1024$ it gives $21.40457$. The cell value is $\Gamma^{J=2}_\infty=21.40452$.
\item \emph{The $1/N$ term.} Near $\theta=0$ the local threshold is $\approx\Gamma_\infty-\kappa\phi^2-c\theta^2$ with $\kappa=0.917$ (phase curvature) and $c=18.16$ ($\theta$ curvature). Along the wall $\theta=i\cdot8/N$ per column. The harmonic-oscillator ground state therefore costs $(8/N)\sqrt{\kappa c}=32.6/N$. Wall curvature adds $+1.325\,\ell$ per unit chord, i.e.\ $+10.6/N$. The prediction $\gamma^\ast(N)\approx\Gamma_\infty-22.0/N$ matches the fitted $21.8/N$. This also explains why Richardson on $N\le64$ (the paper's increments $1.8,0.95$) would overshoot: the $1/N^2$ term is still large there.
\end{itemize}

\section{The leak constant}\label{sec:leak}

\begin{corollary}[Leak constant for every penalty; PROVED]\label{cor:K}
Assume (M0)--(M2), (D), $k\ge4$, and let $\hG\to0$, $h\to0$ with $\mu\ge\mu_0$ arbitrary, in particular with $\gamma\hG\to\infty$.
\begin{enumerate}[label=(\alph*)]
\item On test~P, $K_h:=\norm{\nabla(\ut-u_h)}\mu/(hG)=\norm{\nabla X}/G\to K$, with $|K_h-K|\le C(\hG+(E^0_U)^2)$, $C$ independent of $\mu$, $\gamma$, $\rho$.
The same holds for the leak part $(\mu/h)\delta_R$ of general data.
\item For $\GS$ with general data, $\dn u\cdot\tau\not\equiv0$ on $\Gamma$ and $G>0$, let $K^{\rm gen}_h:=\norm{\nabla(\ut-u_h)}\mu/(hG)$. If $\gamma\hG^{1/2}\to0$ (and the data terms of Corollary~\texttt{u:gen} vanish), then $K^{\rm gen}_h\to K$. If $\gamma\hG^{1/2}\to\infty$, which holds whenever $\gamma\hG\to\infty$, then $K^{\rm gen}_h\to\infty$.
\end{enumerate}
\end{corollary}
\begin{proof}
(a) Corollary~\texttt{u:K} gives $|K_h^2-K^2(1+d_h^2)|\le C(\hG^2+\hG\gamma^{-1/2}(1+\hG^{-1/2}E^0_U))$, and Theorem~\texttt{u:up} gives $d_h\le C(\hG^{1/2}+E^0_U)$. So $|K_h^2-K^2|\le C(\hG+(E^0_U)^2+\hG^{1/2}E^0_U)\le C(\hG+(E^0_U)^2)$, and $K\ge c>0$. For the leak part apply the same results to $\delta_R$ (Corollary~\texttt{hc:cor}(i) decomposition).
(b) The first case is Corollary~\texttt{u:gen}. For the second, Theorem~\texttt{lb:thm} (valid for $\GS$ with every $\gamma\ge\gamma_0$ and every $G$) gives $\norm{\nabla(\ut-u_h)}\ge c_{\rm LB}\kappa_0\norm{\dn u\cdot\tau}_{L^2(\Gamma)}\hG^{3/2}-C\hG^2$. Multiply by $\mu/(hG)=\gamma/(\hG G)$: $K^{\rm gen}_h\ge(c\gamma\hG^{1/2}-C\gamma\hG)/G\to\infty$. Finally $\gamma\hG^{1/2}\ge\gamma\hG$ for $\hG\le1$.
\end{proof}
So the item ``Convergence of the leak constant in the regime $\gamma\hG\to\infty$'' of \texttt{14\_open} is PROVED in the only sense in which it can hold: for the leak itself (test~P, or $\delta_R$). For the total $\GS$ error the normalised constant diverges in that regime. The $\mu=10^3,10^4$ columns of Table~\texttt{tab:P} and of the finite element check in Section~\texttt{hc:sec} are test~P, so they are covered.

\begin{corollary}[Navier--Stokes; PROVED MODULO (R$_S$)]
Under the hypotheses of Theorem~\texttt{n2:up}, $K_{NS,h}:=\norm{\nabla X}/G$ with $X=(\nu\mu/h)\delta_R$ satisfies $|K_{NS,h}-K_{NS}|\le C(\hG^{1/2}+E^0_O)/G$ for every $\gamma\ge\gamma_0$. This is the triangle inequality with Theorem~\texttt{n2:up}. The rate-one Pythagoras argument of \texttt{u:K} used the symmetry of $a_h$ and does not transfer to the Oseen form verbatim.
\end{corollary}

\subsection{$\lim d_h/\hG^{1/2}$ on LP meshes}
Let the family be LP$^+$ (\texttt{11b\_robust} l.162): Lipschitz maps $x\mapsto\mathcal K(x)$ (normalised first-layer shape) and $x\mapsto\ell(x)=\lim|e(x)|/\hG$. For the production meshes, $\ell=\max(\cos^2\theta,\sin^2\theta)$ and $\mathcal K$ is the rectangle lattice of aspect $a(\theta)$ (reduced angle). Let $A(\gamma;\mathcal K)$ be the cell value of \texttt{unifmu2} (\texttt{floquet\_cell.py phi}, $N\to\infty$), i.e.\ $\norm{\nabla W}/(G\ell_{\rm cell}^{1/2})$ for unit data, which is independent of the mode $m$ and of $N$ to $\le0.3\%$ for $N\ge1024$.

\begin{conjecture}[two-scale limit]\label{conj:dh}
For fixed $\mu\ge\mu_0$ (so $\gamma=\mu\hG/h\to\gamma_\infty$) on an LP$^+$ family,
\[
\lim_{h\to0}\frac{d_h^2}{\hG}=D^2:=\frac1{\norm{\nabla U}^2}\int_\Gamma A\bigl(\gamma_\infty\ell(x);\mathcal K(x)\bigr)^2\,\ell(x)\,(p-\pbar)^2\,ds .
\]
\end{conjecture}
\emph{Derivation.} The per-edge energy of $W$ is $A^2c_e^2|e|^2$ with $c_e=(p-\pbar)(m_e^\ast)$. This is the cell normalisation $\norm{\nabla W}^2=A^2G^2\ell$ with $G^2=\sum c_e^2|e|$. Sum over edges with $|e|=\hG\ell$.
\emph{What a proof needs.} (i) Exponential decay of the cell corrector away from the wall (the sawtooth has zero mean on each edge), so that $W\approx\sum_ec_e\,w_{\rm cell}$ with a truncation at depth $O(\log(1/\hG))$ cells. (ii) A transfer (T) of $C^1$ quintic splines from the exact local lattice to the actual LP$^+$ mesh with $O(\hG)$ relative error; the same local-basis theory as (CUT), made quantitative. (iii) Uniform stability, which is available: coercivity for $\gamma\ge\gamma_0$, and Theorem~\texttt{u:up} for the Dirichlet limit. The slow modulation of $c_e$ costs $O(\hG)$ relatively. Status: \lab{CONJECTURE}, PROVED MODULO (T) and (i) along these lines, not written.

\lab{NUMERICAL} (\texttt{leak\_cell.py predict}; $A$ tabulated on ten aspects $a\in[0.75,2.6]$ and interpolated, $\lim h/\hG=3.50$ so $\gamma_\infty=28.6$ at $\mu=100$, local penalties $14.3$--$28.6$, all above the local cell thresholds):
\begin{center}\small
\begin{tabular}{l|cc|c}
\toprule
 & cell $N=1024$ & cell $N=2048$ & fit of production data (\texttt{u:rem})\\
\midrule
$D$, $\mu=10^8$ & 0.5659 & 0.5643 & 0.55\\
$D$, $\mu=100$ & 0.2256 & 0.2256 & 0.225\\
\bottomrule
\end{tabular}
\end{center}
The $\mu=100$ agreement is to $0.3\%$. At $\mu=10^8$ the $3\%$ gap is within the uncertainty of the three-parameter fit on $N\le96$ (its correction factor $1+0.96\hG+3\hG^2$ is still $1.3$ at $N=96$). This supports the conjecture and identifies the limit. The capture fraction $0.41$ of \texttt{u:rem} becomes $D(100)/D(10^8)=0.40$.

\section{Suggested edits (not made)}
\begin{enumerate}
\item \texttt{14\_open} l.26: move ``Convergence of the leak constant in the regime $\gamma\hG\to\infty$'' to \emph{Proved} (Corollary~\ref{cor:K}(a): it is \texttt{u:up}+\texttt{u:K}), adding the divergence for general data (Corollary~\ref{cor:K}(b)). Keep ``existence of $\lim d_h/\hG^{1/2}$'' as numerical, and cite the cell formula (Conjecture~\ref{conj:dh}, agreement $0.3\%$ at $\mu=100$).
\item \texttt{13\_numerics} l.81, last sentence (``in a regime ($\gamma\hG\gg1$) that the proof does not cover (Remark~\texttt{hc:regime})''): now covered by Corollary~\texttt{u:K}.
\item \texttt{14\_open} l.30 and \texttt{13\_numerics} l.114: ``$\gamma^\ast_\infty\approx21.4$'' $\to$ ``$\gamma^\ast_\infty=\gamma_c(\frac34)=21.408$, the periodic half-plane cell value (Theorem~\ref{thm:gam}, PROVED MODULO (CUT), (LOC)), with $\gamma^\ast(N)=21.408-22/N+O(N^{-2})$ observed up to $N=2048$''.
\item \texttt{14\_open} l.36 (``the Navier--Stokes constants $C_R$, $\beta_O$ for any test case''): replace by the numerical values and the resulting $h_0$ (Section~\ref{sec:consts}), stating that the explicit $h_0$ is $\le4\cdot10^{-8}c_\natural$ for test~B at $\nu=1$ and therefore qualitative, while the discrete (L) is observed on every production mesh.
\item \texttt{n2:rem:h0}: add that $U_\infty$ enters through the global maximum of $|\nabla u|$, which here sits at the box corners, far from $\Gamma$.
\end{enumerate}

\section{Files and commands}\label{sec:files}
\texttt{notes/v7/nsconst/}: \texttt{iso\_th.py} (isoparametric Taylor--Hood on the exact $\Omega$); \texttt{consts.py} (\texttt{check N k}, \texttt{CF N k}, \texttt{beta N k nu}, \texttt{CR N k M1,M2,..}, \texttt{Uinf}); \texttt{halfplane.py}; \texttt{discrete\_L.py N nu}; \texttt{cell\_limit.py} (\texttt{phase N a J}, \texttt{table}, \texttt{theta J}, \texttt{mesh N1,N2 J}, \texttt{aspect N}); \texttt{leak\_cell.py} (\texttt{predict N}). Logs: \texttt{run\_beta.log}, \texttt{run\_CR.log}, \texttt{run\_CR16.log}, \texttt{run\_halfplane.log}, \texttt{run\_discreteL.log}, \texttt{run\_table.log}, \texttt{run\_theta.log}, \texttt{run\_mesh.log}, \texttt{run\_predict\_1024.log}, \texttt{run\_predict\_2048.log}. The paper code (\texttt{code/}) and \texttt{notes/v6/unifmu2/floquet\_cell.py} are imported unchanged.

\begin{thebibliography}{9}
\bibitem{Grisvard} P.~Grisvard, \emph{Elliptic Problems in Nonsmooth Domains}, Pitman, 1985. %% VERIFY locators Thm 3.1.1.1 (C^2 part) and Lemma 4.3.1.3 (corners)
\bibitem{LaiSchumaker} M.-J.~Lai, L.~L.~Schumaker, \emph{Spline Functions on Triangulations}, CUP, 2007.
\end{thebibliography}
\end{document}
